\chapter{Laparoscopic radical prostatectomy}
\index{Laparoscopic radical prostatectomy}

\section*{Definition and Overview}
Laparoscopic radical prostatectomy is a minimally invasive surgical procedure performed to remove the entire prostate gland, seminal vesicles, and adjacent tissue, typically as a primary curative treatment for localised prostate cancer. Unlike traditional open radical retropubic prostatectomy, laparoscopic radical prostatectomy is performed using small keyhole incisions in the abdomen through which a laparoscope (a camera-tipped tube) and specialised surgical instruments are inserted. This approach provides the surgeon with a magnified, high-definition three-dimensional view of the pelvic anatomy, facilitating precise dissection of the prostate while sparing surrounding critical structures, such as the cavernous nerves responsible for erectile function and the external urethral sphincter responsible for urinary continence. The procedure may be performed manually or with robotic assistance (robotic-assisted laparoscopic prostatectomy), both aiming to achieve oncological cure while minimising postoperative morbidity, reducing blood loss, and accelerating patient recovery.

\section*{Epidemiology}
Prostate cancer is the most commonly diagnosed non-cutaneous cancer in men globally, with particularly high incidence rates in developed nations, including some countries and the United Kingdom. In many countries, it represents a major public health concern, with thousands of cases diagnosed annually, predominantly in men aged 60 years and older. Laparoscopic and robotic-assisted radical prostatectomies have largely superseded open surgery as the preferred surgical approach in high-income countries due to reduced blood loss, decreased transfusion rates, and shorter hospital stays. Risk factors for prostate cancer itself include advancing age, a family history of the disease, specific genetic mutations (such as BRCA1/2), and African ancestry. While radical prostatectomy is primarily offered to patients with localised (T1 or T2) or selected locally advanced (T3) disease who have a life expectancy of at least 10 years, the choice of the laparoscopic approach varies depending on surgeon expertise, hospital resources, and patient-specific anatomical factors.

\section*{Aetiology and Pathophysiology}
Radical prostatectomy is indicated for the treatment of prostate cancer, which arises from malignant transformation of epithelial cells in the prostate gland, most commonly within the peripheral zone. The pathology is characterised by uncontrolled cellular proliferation driven by androgenic hormones, leading to tissue invasion and potential metastatic spread.
\begin{itemize}
    \item \textbf{Oncological indication:} Primarily indicated for clinically localised prostate cancer (Gleason score 6 to 10, or ISUP grade groups 1 to 5) where the cancer is confined to the prostate gland, with the goal of complete tumour eradication.
    \item \textbf{Anatomical disruption:} The surgical removal of the prostate involves transecting the urethra at the apex of the prostate and the bladder neck at the base, followed by reconstruction via a vesicourethral anastomosis.
    \item \textbf{Neurovascular injury risk:} The cavernous nerves run in close proximity to the posterolateral aspect of the prostate. Surgical dissection can cause temporary or permanent damage due to traction, thermal injury from cautery, or direct transection, leading to postoperative erectile dysfunction.
    \item \textbf{Sphincteric impairment:} The external urethral sphincter lies just distal to the prostatic apex. Avoidance of damage to this structure and preservation of maximal functional urethral length are critical to preventing postoperative urinary incontinence.
\end{itemize}

\section*{Clinical Features}
Patients undergoing laparoscopic radical prostatectomy typically present with features of their underlying prostate cancer or may be completely asymptomatic, with the disease detected via routine screening.
\begin{itemize}
    \item Asymptomatic presentation is common, where the cancer is detected solely due to an elevated prostate-specific antigen (PSA) level or abnormal digital rectal examination (DRE).
    \item Lower urinary tract symptoms (LUTS), including urinary frequency, urgency, nocturia, or a weak urinary stream, may occur if the tumour is large or coexists with benign prostatic hyperplasia.
    \item Postoperative features include abdominal discomfort at the keyhole incision sites, transient haematuria, and the presence of an indwelling urinary catheter to allow the vesicourethral anastomosis to heal.
    \item Early postoperative urinary leakage or incontinence upon catheter removal is common, presenting as stress urinary incontinence.
    \item Postoperative erectile dysfunction is a frequent feature, characterised by the inability to achieve or maintain erections sufficient for penetration, which may gradually improve over 12 to 24 months.
\end{itemize}

\section*{Differential Diagnosis}
While laparoscopic radical prostatectomy is a treatment rather than a primary diagnosis, the clinical presentation of prostate cancer must be distinguished from other prostatic and urological conditions:
\begin{itemize}
    \item \textbf{Benign prostatic hyperplasia:} A non-malignant enlargement of the prostate transition zone that causes similar lower urinary tract symptoms and mild PSA elevation.
    \item \textbf{Prostatitis:} Acute or chronic inflammation of the prostate, which can cause significant pelvic pain, dysuria, urinary frequency, and marked, transient elevations in PSA.
    \item \textbf{Prostate stones or infarction:} Benign prostatic conditions that can occasionally cause lower urinary tract symptoms, haematuria, or localised prostatic tenderness.
    \item \textbf{Bladder neck contracture or urethral stricture:} Conditions that cause urinary obstruction and may mimic or complicate postoperative voiding dysfunction.
\end{itemize}

\section*{When to Seek Medical Care}
Patients who have undergone laparoscopic radical prostatectomy receive detailed discharge instructions. They should contact their urology team or seek prompt medical attention if they experience issues with their urinary catheter, signs of infection, or calf swelling.

\textbf{Contact your local emergency services for:}
\begin{itemize}
    \item Sudden, severe shortness of breath or chest pain (suggestive of pulmonary embolism).
    \item Heavy, continuous bleeding from the wound sites or passing large blood clots in the urine.
    \item Signs of shock, such as severe dizziness, confusion, cold or clammy skin, or loss of consciousness.
\end{itemize}

Other non-emergency reasons to contact the doctor or surgical team include:
\begin{itemize}
    \item The urinary catheter stops draining urine or falls out prematurely.
    \item Fever above 38 degrees Celsius, shaking chills, or worsening abdominal pain.
    \item Worsening redness, swelling, warmth, or purulent drainage from the keyhole wounds.
    \item Severe swelling or pain in one calf or leg.
\end{itemize}

\section*{Diagnosis and Investigations}
The diagnosis of prostate cancer and preparation for laparoscopic radical prostatectomy involve several laboratory, imaging, and histopathological tests.
\begin{itemize}
    \item \textbf{Prostate-specific antigen (PSA) test:} A blood test measuring a glycoprotein produced by prostatic tissue; elevated levels or a rapid rise over time raise suspicion of prostate cancer.
    \item \textbf{Digital rectal examination (DRE):} A physical examination where the clinician inserts a gloved finger into the rectum to assess the size, shape, and consistency of the prostate, looking for hard nodules or asymmetry.
    \item \textbf{Multiparametric MRI (mpMRI):} A specialised imaging study of the prostate used to identify suspicious areas (scored via the PI-RADS system) and assist in targeting biopsies and planning nerve-sparing surgery.
    \item \textbf{Transrectal or transperineal prostate biopsy:} The definitive diagnostic test, involving ultrasound-guided tissue sampling of the prostate to determine the Gleason score and ISUP grade group.
    \item \textbf{Staging investigations:} Bone scans or PSMA PET-CT scans are performed in intermediate-risk or high-risk patients to rule out distant metastases before surgery.
    \item \textbf{Pre-operative assessment:} Routine blood tests, including full blood count, electrolytes, renal function tests, coagulation profile, and an electrocardiogram (ECG) to confirm fitness for general anaesthesia.
\end{itemize}

\section*{Management}
Laparoscopic radical prostatectomy is a complex surgical procedure that requires comprehensive perioperative and postoperative management.
\begin{itemize}
    \item \textbf{Surgical procedure:} Performed under general anaesthesia, involving 5 to 6 small abdominal incisions, dissection of the prostate, seminal vesicles, and often pelvic lymph nodes, followed by reconstruction of the urinary tract (vesicourethral anastomosis) over a urinary catheter.
    \item \textbf{Thromboembolism prophylaxis:} Use of compression stockings, sequential compression devices, and pharmacological anticoagulants (such as low-molecular-weight heparin) to prevent deep vein thrombosis and pulmonary embolism.
    \item \textbf{Postoperative pain management:} Multimodal analgesia utilising paracetamol, non-steroidal anti-inflammatory drugs (NSAIDs), and local anaesthetic infiltration, minimising the use of opioids to reduce the risk of postoperative ileus.
    \item \textbf{Urinary catheter management:} The indwelling urethral catheter is typically left in place for 7 to 14 days postoperatively to support the healing of the new bladder-urethra connection; patients are trained in catheter care prior to discharge.
    \item \textbf{Pelvic floor muscle training:} Pelvic floor exercises (Kegels) are initiated pre-operatively and resumed postoperatively after catheter removal to strengthen the external sphincter and accelerate recovery of urinary continence.
    \item \textbf{Penile rehabilitation:} Postoperative use of phosphodiesterase-5 (PDE5) inhibitors (such as sildenafil or tadalafil), vacuum erection devices, or intracavernosal injections to maintain penile tissue oxygenation and promote recovery of erectile function.
    \item \textbf{Oncological follow-up:} Monitoring of blood PSA levels, typically starting 6 to 8 weeks post-surgery and repeated at regular intervals; an undetectable PSA level (typically <0.1 or <0.2 ng/mL) indicates successful surgical removal.
\end{itemize}

\section*{Complications}
\begin{itemize}
    \item \textbf{Urinary incontinence:} Stress urinary incontinence is very common immediately after catheter removal; while most patients regain continence within a year, a small percentage experience persistent leakage requiring pads or surgical intervention (e.g., male sling or artificial urinary sphincter).
    \item \textbf{Erectile dysfunction:} A common long-term complication, even with bilateral nerve-sparing techniques, due to neuropraxia (nerve stretching) or direct injury; recovery can take up to 2 years and may be incomplete.
    \item \textbf{Anastomotic leak or stricture:} Leakage of urine at the vesicourethral junction can lead to pelvic fluid collections, while scarring can cause bladder neck contracture, resulting in voiding difficulties.
    \item \textbf{Bleeding and transfusion:} Intraoperative haemorrhage may occur, occasionally requiring blood transfusion or conversion to open surgery.
    \item \textbf{Infection:} Surgical site infection, pelvic abscess, or urinary tract infection associated with catheter use.
    \item \textbf{Thromboembolic events:} Deep vein thrombosis or pulmonary embolism due to pelvic surgery and immobility.
    \item \textbf{Adjacent organ injury:} Rare risk of rectal injury or obturator nerve damage during pelvic dissection.
\end{itemize}

\section*{Prognosis and Outcomes}
The prognosis after laparoscopic radical prostatectomy is generally excellent, particularly for patients with localised, low-to-intermediate grade prostate cancer. The 10-year cancer-specific survival rate exceeds 90--95\% for local disease. Long-term oncological success depends on the pathological stage, Gleason score, and surgical margin status determined from the excised prostate specimen. If the cancer recurs (indicated by a rising PSA, known as biochemical recurrence), salvage therapies such as radiation therapy or androgen deprivation therapy remain highly effective. The functional prognosis for urinary continence is favourable, with over 85--90\% of patients regaining satisfactory control within 12 months. Recovery of erectile function is highly variable and depends on pre-operative potency, patient age, and whether a unilateral or bilateral nerve-sparing procedure was performed.

\section*{Prevention and Self-Care}
\begin{itemize}
    \item Perform pelvic floor exercises exactly as instructed by your physiotherapist or urologist, starting before the surgery and continuing after catheter removal.
    \item Care for the urinary catheter by washing the insertion site daily, keeping the drainage bag below the level of the bladder, and drinking plenty of fluids to prevent blockage.
    \item Avoid heavy lifting (greater than 5 kg) and strenuous physical activity for at least 4 to 6 weeks post-surgery to allow the abdominal muscles and anastomosis to heal.
    \item Maintain a diet high in fibre and use stool softeners as prescribed to prevent straining during bowel movements, which can put pressure on the surgical site.
    \item Walk regularly at a gentle pace to promote circulation, improve energy levels, and reduce the risk of blood clots.
    \item Attend all scheduled follow-up appointments for wound checks, catheter removal, and PSA blood tests.
\end{itemize}

\section*{Sources and Further Reading}
The following reputable organisations provide further information and support for laparoscopic radical prostatectomy and prostate cancer management:
\begin{itemize}
    \item Prostate Cancer Foundation of Australia. \textit{Understanding surgery: radical prostatectomy guide.} https://www.pcfa.org.au/
    \item Cancer Council Australia. \textit{Prostate cancer diagnosis, treatment, and recovery.} https://www.cancer.org.au/
    \item NHS. \textit{Prostate cancer treatment options and radical prostatectomy.} https://www.nhs.uk/conditions/prostate-cancer/treatment/
    \item Mayo Clinic. \textit{Radical prostatectomy: what you can expect.} https://www.mayoclinic.org/tests-procedures/prostatectomy/
    \item European Association of Urology. \textit{Patient information on radical prostatectomy.} https://patients.uroweb.org/
\end{itemize}
